\subsection{The exact all-shell no-hit obstruction}
\label{sec:all-shell-no-hit}

Fix a prime \(p=12h+1\), \(h\geq1\), and retain the original shell interval
\[
 A_p=\{3h+1,\ldots,9h\}.
\]
For each \(a\in A_p\) put \(R_a=4a-p\). The unit lemma already
proved above gives \(\gcd(a,R_a)=\gcd(p,R_a)=1\). Write the actual
prime factorization, with no suppressed factors,
\[
 a=\prod_{\ell\mid a}\ell^{v_\ell(a)},\qquad
 {\cal B}_a=\prod_{\ell\mid a}\{0,1,\ldots,2v_\ell(a)\}.
\]
The map
\[
 \nu_a:{\cal B}_a\longrightarrow\operatorname{Div}(a^2),\qquad
 (e_\ell)_{\ell\mid a}\longmapsto U_a(e):=\prod_{\ell\mid a}\ell^{e_\ell}
 \tag{1}
\]
is bijective. Its inverse is the full valuation vector
\(e_\ell=v_\ell(u)\); in particular, no exponent interval is
replaced by the subgroup it generates.

Let \(U(R_a)=(\mathbb Z/R_a\mathbb Z)^\times\), written
multiplicatively, and let \(\mathbb N[U(R_a)]\) be its integral
semiring of finite formal sums. Since every \(\ell\mid a\) is a
unit modulo \(R_a\), define the finite unit-group packet
\[
 F_{a,R_a}:=\prod_{\ell\mid a}
   \left(\sum_{e=0}^{2v_\ell(a)}[\ell]^e\right)
 =\sum_{g\in U(R_a)}c_{a,R_a}(g)[g].
 \tag{2}
\]
The coordinate reduction map is the typed map
\[
 \pi_{a,R_a}:\operatorname{Div}(a^2)\longrightarrow U(R_a),
 \qquad u\longmapsto [u]_{R_a},
\]
and the exponent-to-unit map is $\pi_{a,R_a}\circ\nu_a$.
Expanding (2) records one and only one original exponent at every
prime. Thus, for every \(g\in U(R_a)\),
\[
 c_{a,R_a}(g)=
 \#\{e\in{\cal B}_a:U_a(e)\equiv g\pmod {R_a}\}. \tag{3}
\]
The two original channel sets are
\[
 E_a=\{u\mid a^2:R_a\mid4u+1\},\qquad
 M_a=\{u\mid a^2:R_a\mid u+a\}. \tag{4}
\]
Their exact coefficient form is
\[
 |E_a|=c_{a,R_a}(-4^{-1}),\qquad
 |M_a|=c_{a,R_a}(-a). \tag{5}
\]
Indeed, \(4\) and \(a\) are units modulo \(R_a\), so (4) is exactly
the pair of residue equations \(u\equiv-4^{-1}\) and
\(u\equiv-a\), and (3) counts the complete valuation fibres.

For later reconstruction, define the direct residual maps
\[
 \rho^E_a,\rho^M_a:{\cal B}_a\longrightarrow\mathbb Z/R_a\mathbb Z,
 \qquad
 \rho^E_a(e)=4U_a(e)+1\pmod {R_a},\qquad
 \rho^M_a(e)=U_a(e)+a\pmod {R_a}. \tag{6}
\]
The fibres of the zero element of these maps are, by (1),
\[
 (\rho^E_a)^{-1}(0)=\nu_a^{-1}(E_a),\qquad
 (\rho^M_a)^{-1}(0)=\nu_a^{-1}(M_a), \tag{7}
\]
and the inverse on either fibre is the coordinatewise valuation
map \(u\mapsto(v_\ell(u))_{\ell\mid a}\). Equations (1)--(7) are
coordinate-level maps between the exponent box, the divisor set,
and the finite unit group; they retain every prime and every
valuation bound.

\begin{theorem}[Shellwise no-hit criterion and integer return]
\label{thm:shellwise-no-hit}
For \(a\in A_p\), the shell has no original witness if and only if
\[
 c_{a,R_a}(-4^{-1})=0
 \quad\hbox{and}\quad
 c_{a,R_a}(-a)=0. \tag{8}
\]
The complete set of ordered witnesses with first denominator \(a\) is
\[
 \mathcal V_{a,p}^{\rm ord}
 =\{(a,y,z):y,z\in\mathbb Z_{>0},\quad 1/a+1/y+1/z=4/p\}.
\]
Its exact tagged coordinate domain is the disjoint union
\[
 \mathcal T_{a,p}
 =\{(0,u,\sigma):u\in E_a,\ \sigma\in\{0,1\}\}
 \ \coprod\ \{(1,u,0):u\in M_a\}.
\]
For a point \((\varepsilon,u,\sigma)\) in this domain, let
\[
 \begin{aligned}
 g&=\gcd(a,u),& A&=a/g,\\
 B&=u/g,&D&=g^2/u,\qquad
 C=\frac{A+p^{1-\varepsilon}B}{R_a}.
 \end{aligned}\tag{9}
\]
Then \(A,B,D,C\) are positive integers, \(a=ABD\), \(u=B^2D\),
and \(\gcd(A,B)=1\). Preserve the original chart values
\[
 y_0=p^\varepsilon ACD,\qquad z_0=pBCD.
\]
The map
\[
 \begin{gathered}
 \Phi_{a,p}:\mathcal T_{a,p}\longrightarrow\mathcal V_{a,p}^{\rm ord},\\
 \Phi_{a,p}(\varepsilon,u,\sigma)=
 \begin{cases}
  (a,y_0,z_0),&\sigma=0,\\
  (a,z_0,y_0),&\sigma=1
 \end{cases}
 \end{gathered} \tag{10}
\]
is a bijection. For its full inverse, retain
\(S_a=pa\), \(d_y=R_a y-S_a\), \(d_z=R_a z-S_a\), and put
\(i=v_p(d_y)\). Then \(d_y d_z=S_a^2\), \(i\in\{0,1,2\}\), and
\[
 \Phi_{a,p}^{-1}(a,y,z)=
 \begin{cases}
  (0,a^2/d_y,0),&i=0,\\
  (1,pa^2/d_y,0),&i=1,\\
  (0,d_y/p^2,1),&i=2.
 \end{cases} \tag{10a}
\]
Equations (9) and the valuation inverse in (1) then recover every
chart coordinate and every original exponent. The ordered cardinality is
\(2|E_a|+|M_a|\).

\begin{proof}
For \(u\mid a^2\), write \(v=v_\ell(a)\), \(e=v_\ell(u)\) at
each original prime \(\ell\mid a\). Then \(0\le e\le2v\), and
the exponents of \(A,B,D\) in (9) are respectively
\[
 v-\min(v,e),\qquad e-\min(v,e),\qquad 2\min(v,e)-e.
\]
All are nonnegative, the first two cannot both be positive, and
their sums for \(ABD\) and \(B^2D\) are \(v\) and \(e\).
This proves integrality, both product identities, and coprimality.
Conversely these chart coordinates recover \(u=B^2D\) and
\(g=BD\), hence every original exponent. The gate in (4) is
equivalent to \(4u+p^\varepsilon\equiv0\pmod{R_a}\). Indeed for
\(\varepsilon=1\) use \(p\equiv4a\) and the unit \(4\).
Keeping all factors, its chart form follows from
\[
 A(4u+p^\varepsilon)
 \equiv pB+p^\varepsilon A
 =p^\varepsilon(A+p^{1-\varepsilon}B)\pmod{R_a}.
\]
The multipliers \(A\) and \(p\) are units modulo \(R_a\), so the
equivalence holds in both directions and \(C\) is a positive
integer. Substitution of the unswapped chart, with common
denominator \(pABCD\), gives the exact integer identity
\[
 \frac{pABCD}{a}+\frac{pABCD}{p^\varepsilon ACD}
 +\frac{pABCD}{pBCD}
 =pC+p^{1-\varepsilon}B+A
 =4ABCD,
\]
because \((4ABD-p)C=R_aC=A+p^{1-\varepsilon}B\).
Thus (10) is a positive integer witness in either retained order.
The original residuals in the unswapped chart are exactly
\[
 \begin{aligned}
 R_a y_0-S_a
  &=p^\varepsilon A D(A+p^{1-\varepsilon}B)-pABD\\
  &=p^\varepsilon A^2D=p^\varepsilon a^2/u,\\
 R_a z_0-S_a
  &=pBD(A+p^{1-\varepsilon}B)-pABD\\
  &=p^{2-\varepsilon}B^2D=p^{2-\varepsilon}u.
 \end{aligned} \tag{10b}
\]
Swapping \(y_0,z_0\) swaps these two residuals.

For completeness, the typed residual domain is
\[
 \mathcal D_{a,p}
 =\{d\in\operatorname{Div}(S_a^2):R_a\mid d+S_a\}.
\]
The map \(\mathcal V_{a,p}^{\rm ord}\to\mathcal D_{a,p}\) sends
\((a,y,z)\) to \(d_y=R_a y-S_a\). The original equation gives
\(R_a yz=S_a(y+z)\), so \(d_y=S_a y/z>0\) and
\(d_y d_z=S_a^2\). Its inverse is
\[
 d\longmapsto
 \left(a,\frac{S_a+d}{R_a},\frac{S_a+S_a^2/d}{R_a}\right).
 \tag{10c}
\]
The first quotient is integral by the gate. Since every divisor
of \(S_a^2\) is a unit modulo \(R_a\), multiplication of the second
numerator by \(d\) gives \(S_a(d+S_a)\equiv0\); therefore the second
quotient is also integral. Both are positive, and expanding the
residual product gives \(R_a yz=S_a(y+z)\), hence the required
reciprocal equation. The formulas recover both residuals and both
denominators, so both inverse compositions are the identity.

As \(p\nmid a\), every \(d\in\mathcal D_{a,p}\) has a unique form
\(d=p^i v\) with \(i\in\{0,1,2\}\), \(v\mid a^2\).
We now compute all three inverse branches without deleting one of them.
For \(i=0\), put \(u=a^2/d\). Multiplication of
\(d+pa\equiv0\) by the unit \(u\), followed by division by the
unit \(a\), gives \(a+pu\equiv0\), which is exactly
\(4u+1\equiv0\) after using \(p\equiv4a\).
Thus \(u\in E_a\), and (10b) recovers \(d_y=a^2/u=d\)
with \((\varepsilon,\sigma)=(0,0)\).
For \(i=2\), put \(u=d/p^2=v\). Division of the gate by the unit
\(p\) gives \(pu+a\equiv0\), again equivalent to the exterior
gate. The swapped chart has \(d_y=p^2u=d\), with
\((\varepsilon,\sigma)=(0,1)\).
For \(i=1\), write \(d=pv\) and put \(u=a^2/v=pa^2/d\).
The gate is \(v+a\equiv0\). Multiplication by the unit \(u\)
and division by the unit \(a\) gives \(a+u\equiv0\).
Thus \(u\in M_a\), and (10b) recovers \(d_y=pa^2/u=d\)
with \((\varepsilon,\sigma)=(1,0)\).
Each modular multiplication and division used a unit, so every
gate equivalence has its converse. These branches prove that
(10a) maps every ordered witness into the asserted tagged domain
and that (10) recovers the original \(d_y\), hence the original
ordered triple by (10c). Conversely (10b) has \(p\)-adic exponent
respectively \(0,2,1\) in the exterior unswapped, exterior swapped,
and middle branches. Formula (10a) therefore recovers the original
\(\varepsilon,u,\sigma\) in each case. Every fibre is a singleton.
The domain has cardinality \(2|E_a|+|M_a|\). It is empty exactly
when both coefficients in (5) vanish, proving (8).
\end{proof}
\end{theorem}

The exact map corresponding to exchange of the last two original
denominators is the involution
\(\iota_{a,p}:\mathcal T_{a,p}\to\mathcal T_{a,p}\) given by
\[
 (0,u,\sigma)\longmapsto(0,u,1-\sigma),\qquad
 (1,u,0)\longmapsto(1,a^2/u,0). \tag{10d}
\]
For the middle branch \(R_a\mid u+a\) implies
\(R_a\mid a^2/u+a\), by multiplication by the unit
\([a]_{R_a}[u]_{R_a}^{-1}\in U(R_a)\).
Its chart changes \((A,B,C,D)\) to \((B,A,C,D)\): indeed
\(a^2/u=A^2D\), its gcd with \(a=ABD\) is \(AD\), and
\(C=(A+B)/R_a\) is unchanged. This proves the exchange in
(10d) directly in the original coordinates. There is no extra
middle swap bit, because it would duplicate these same ordered
witnesses. Define the unordered codomain and the two order-forgetting
maps explicitly by
\[
 \mathcal V_{a,p}^{\rm unord}
 =\{(a,\{y,z\}):(a,y,z)\in\mathcal V_{a,p}^{\rm ord}\},
\]
\[
 \begin{aligned}
 q_{\mathcal V}:\mathcal V_{a,p}^{\rm ord}&\longrightarrow
    \mathcal V_{a,p}^{\rm unord},&
 (a,y,z)&\longmapsto(a,\{y,z\}),\\
 q_{\mathcal T}:\mathcal T_{a,p}&\longrightarrow
    \mathcal V_{a,p}^{\rm unord},&
 q_{\mathcal T}&=q_{\mathcal V}\circ\Phi_{a,p}.
 \end{aligned}
\]
The exact fibres of \(q_{\mathcal T}\) are
\[
 \{(0,u,0),(0,u,1)\},\qquad
 \{(1,u,0),(1,a^2/u,0)\}.
\]
The corresponding \(q_{\mathcal V}\)-fibres are
\(\{(a,y,z),(a,z,y)\}\), by the proved bijection \(\Phi_{a,p}\).
These fibres all have two elements: a fixed middle point would have
\(u=a\) and \(R_a\mid2a\), forcing \(R_a\mid2\), impossible.
Equivalently a witness with \(y=z\) would have \(d_y=S_a\),
forcing \(R_a\mid2S_a\), again impossible by the unit lemma.

Define the shell count and global obstruction without identifying
different shell labels:
\[
 q_{a,p}=|E_a|+\frac{|M_a|}{2},\qquad
 Q_p=\sum_{a\in A_p}q_{a,p},\qquad
 {\mathcal O}_p=\prod_{a\in A_p}
 {\bf1}\!\left[c_{a,R_a}(-4^{-1})=0\right]
 {\bf1}\!\left[c_{a,R_a}(-a)=0\right]. \tag{11}
\]
The middle involution \(u\mapsto a^2/u\) has no fixed point on
\(M_a\): a fixed point would be \(u=a\), forcing
\(R_a\mid2a\), hence \(R_a\mid2\), impossible since \(R_a\ge3\).
Thus \(|M_a|\) is even. The proved two-element fibres show that
\(q_{a,p}\) is exactly the number of unordered pairs underlying
\(\mathcal V_{a,p}^{\rm ord}\), and
\(2q_{a,p}=|\mathcal V_{a,p}^{\rm ord}|\).

\begin{theorem}[Necessary-and-sufficient global obstruction]
\label{thm:global-no-hit}
For every prime \(p=12h+1\),
\[
 {\mathcal O}_p=1
 \Longleftrightarrow
 (\forall a\in A_p)\ E_a=M_a=\varnothing
 \Longleftrightarrow Q_p=0
 \Longleftrightarrow
 \text{\(p\) has no original Erd\H{o}s--Straus witness}. \tag{12}
\]
Consequently a hypothetical counterexample must satisfy the finite,
computable condition
\[
 \forall a\in\{3h+1,\ldots,9h\}:\quad
 [\,-4^{-1}\,]\notin\operatorname{supp}(F_{a,R_a})
 \quad\text{and}\quad
 [\,-a\,]\notin\operatorname{supp}(F_{a,R_a}), \tag{13}
\]
with every support computed from the unabridged boxes \({\cal B}_a\).

\begin{proof}
Each factor in (11) is \(0\) or \(1\), and equals \(1\) exactly
when the corresponding complete fibre is empty by (5)--(7).
Therefore the product is \(1\) exactly when every shell has both
empty channels. The summands \(q_{a,p}\) are nonnegative and
vanish exactly under the same condition, so their finite sum
vanishes exactly then. The full primewise biconditional already
proved above is \(\operatorname{ES}(p)\Longleftrightarrow Q_p>0\);
taking its negation yields the last equivalence. Since (2) is a
finite product, (13) is a finite prime-factor algorithm and not an
occupancy or density assertion.
\end{proof}
\end{theorem}

\begin{corollary}[Residual-three and residual-seven tests]
\label{cor:global-r3-r7}
For \(a=3h+1\), the coefficient in (5) is nonzero exactly when
\(a\) has a prime factor \(\ell\equiv2\pmod3\); this is the
complete residual-three sieve. For \(a=3h+2\), the two
coefficients in (5) are given by the six-class packet and the
three explicit predicates \(n_3,n_{24}\) in
Theorems~\ref{thm:r3-factor-sieve}--\ref{thm:r7-factor-sieve}.
Thus a global no-hit \(p\) must pass both exact local rejection tests,
together with (13) for every remaining shell. These conditions
are necessary and sufficient and retain the complete factor and
shell data.
\end{corollary}

\paragraph{A complete original-order fixture.}
At \(p=13\), \(h=1\), \(a=4\), one has \(R_a=3\), \(S_a=52\),
\(a^2=16\), and \(E_a=M_a=\{2,8\}\). These are exactly the
divisors in \(\{1,2,4,8,16\}\) congruent to \(2\pmod3\).
All six tagged points and their original residuals and denominators are
\[
\begin{array}{c|r|r|r|r}
 (\varepsilon,u,\sigma)&d_y&d_z&y&z\\\hline
 (0,2,0)&8&338&20&130\\
 (0,2,1)&338&8&130&20\\
 (0,8,0)&2&1352&18&468\\
 (0,8,1)&1352&2&468&18\\
 (1,2,0)&104&26&52&26\\
 (1,8,0)&26&104&26&52
\end{array}
\]
Here \(S_a^2=2704=2^4\cdot13^2\). Its passing divisors are
exactly \(2,8,26,104,338,1352\): among
\(2^e13^i\), \(0\le e\le4\), \(0\le i\le2\), the gate
\(d_y+52\equiv0\pmod3\) holds exactly for \(e=1,3\).
Formula (10c) gives precisely the six rows above. The shell has
three unordered pairs and six ordered witnesses; restricting the
chart to \(\sigma=0\) alone would omit the two exterior rows
with \(v_{13}(d_y)=2\).

\paragraph{A checked full-fibre warning.}
At $p=373=12\cdot31+1$ and $a=95=5\cdot19$, one has
$R_a=7$ and
$F_{95,7}=(1+[5]+[4])^2$ in $\mathbb N[U(7)]$ (the two
factors are separately retained, even though they have the same
residue).  The complete exterior fibre is
$E_{95}=\{5,19\}$ and the middle fibre is empty, since the
exponent sums $i+j$ with $0\le i,j\le2$ produce residues
$1,5,4,6,2$ and never $-95\equiv3\pmod7$.  The two retained
exponent vectors reconstruct $(A,B,C,D,\varepsilon)=(19,1,56,5,0)$
and $(5,1,54,19,0)$, so every hit in this shell has $A>2$.
This is a direct finite-fibre check and does not replace the full
all-shell criterion by an $A\le2$ cutoff.

\paragraph{Checked finite algorithm.}
For each input prime \(p=12h+1\), factor every \(a\in A_p\), enumerate
the literal exponent box \({\cal B}_a\), accumulate residues in
\(U(R_a)\), and test the two target coefficients. The checker
\texttt{check\_nohit.py} implements (1)--(13), compares each
coefficient with direct divisor enumeration, verifies (9)--(10d),
compares the complete ordered tagged image with direct enumeration
of every passing divisor of the original \(S_a^2\), checks both
inverse compositions and every order-forgetting fibre, and records
the first shell having a hit. It does not claim that
every shell is occupied, that a prescribed residue is hit by a
prime factor, or that infinitely many survivors exist.
